\chapter{Flask Web Application}\label{ch:flask}

\section{Inference module}
\file{src/visionguard/inference.py} defines \texttt{Predictor}. Its constructor runs once per
process: it loads \file{champion.keras} with \texttt{keras.models.load\_model}, reads
\file{champion.json} (image size, resize mode, threshold, version), and runs one dummy prediction
to warm up the TensorFlow graph. \texttt{predict\_bytes(data)} decodes the bytes, resizes them,
predicts $p(\text{AI})$, applies the threshold and returns label, confidence (probability of the
predicted label), $p(\text{AI})$, threshold, model name/version and latency.

\textbf{Why not load per request.} Loading deserialises millions of weights and builds the graph:
it takes seconds and hundreds of MB. Per request, every user would wait for it, memory would
spike with concurrent requests, and the first predict of a fresh model is the slowest one. Loading
once and keeping the model in memory makes each prediction tens of milliseconds.

\textbf{Why training and serving preprocessing must match.} The model only learned the mapping
from \emph{the inputs it saw}. If serving resized with a different interpolation, kept the alpha
channel, scaled to a different range or preserved the aspect ratio when training did not, the
inputs would come from a different distribution and predictions would silently degrade, with no
error raised (\emph{training/serving skew}). I removed the possibility by construction: training,
audit and serving call the same \texttt{decode\_and\_resize} function, the image size and resize
mode travel with the model in \file{champion.json}, and backbone-specific scaling is a layer
inside the saved model.

\shot{s4_1_console}{\file{scripts/infer_cli.py}: model loaded once, then two known test images (one AI, one real).}

\section{Routes and validation}
\begin{tabularx}{\linewidth}{>{\ttfamily}l l X}
\toprule
Route & Method & Purpose \\
\midrule
/ & GET & upload page (drag \& drop, client-side preview, last 5 checks) \\
/predict & POST & validate, predict, write DB row, redirect to the result (Post/Redirect/Get) \\
/result/<id> & GET & verdict, confidence, $p(\text{AI})$ vs threshold, caveat, audit record \\
/history & GET & paginated database trail, filter by verdict, show rejected uploads \\
/api/history & GET & the same trail as JSON \\
/health & GET & JSON health check: DB reachable, model backend reachable, model version \\
\bottomrule
\end{tabularx}

\textbf{Validation} (\file{src/visionguard/validation.py}), cheapest check first:
\begin{enumerate}
\item No file $\rightarrow$ 400.
\item Extension not in \{jpg, jpeg, png, webp, bmp, gif\} $\rightarrow$ \textbf{415 Unsupported file type}.
\item Size: Flask's \texttt{MAX\_CONTENT\_LENGTH} rejects oversized bodies before my code runs,
and the stream is read with a hard limit of 5\,MB + 1 byte $\rightarrow$ \textbf{413 File too large}.
\item Full decode with Pillow (verify + load) $\rightarrow$ \textbf{422 Not a readable image} for
renamed non-images or truncated files. Pillow's decompression-bomb limit is also set.
\end{enumerate}
File names go through \texttt{werkzeug.utils.secure\_filename} (e.g.\
\texttt{../../etc/passwd.png} becomes \texttt{etc\_passwd.png}) and are only stored as text.
The upload is read into memory and \textbf{never written to disk}: only its SHA-256 hash, size
and dimensions are kept. Every error renders a friendly page with the status code, a plain
explanation and a way back; unexpected exceptions are logged with a traceback but the user only
sees a generic 500 page.

\shot[0.8]{s4_2_wrong_type}{Failure case 1: wrong file type (415).}
\shot[0.8]{s4_2_too_large}{Failure case 2: file over the size limit (413).}
\shot[0.8]{s4_2_corrupt}{Failure case 3: a text file renamed to \texttt{.jpg} (422).}

\section{HTML and CSS}
Templates use Jinja inheritance (\file{base.html} holds the header, navigation with the History
menu, and the footer showing the model version and inference mode). The result page shows the
verdict in the class colour, a confidence bar, a scale with the decision threshold and the
model's $p(\text{AI})$, a ``low certainty'' badge under 65\%, and an honest caveat box: the result
is a statistical estimate, weakest on screenshots, heavy compression and newer generators, tuned to
avoid false accusations, and never sufficient on its own. The stylesheet is plain CSS with a dark
mode and a mobile layout.

\shot{s4_3_upload}{Upload page.}
\shot{s4_3_result}{Result page with verdict, confidence, threshold scale and caveat.}
\shot{s4_3_history}{History menu: the database trail.}

\begin{wrongbox}
With the threshold chosen on validation data below 0.5 (e.g.\ 0.47), an image with
$p(\text{AI})=0.49$ is labelled AI-generated with a ``confidence'' of 49\%, which looks
contradictory. The number is honest (it is the model's probability for that label), so instead of
hiding it the page now shows the threshold scale and a ``low certainty'' badge.
\end{wrongbox}

\section{Configuration and production server}
All settings are environment variables read in \file{app/config.py}: \texttt{MODEL\_PATH},
\texttt{THRESHOLD} (optional override), \texttt{DATABASE\_URL}, \texttt{INFERENCE\_MODE}
(\texttt{local}/\texttt{remote}), \texttt{MODEL\_API\_URL}, \texttt{MODEL\_API\_KEY},
\texttt{MAX\_UPLOAD\_MB}, \texttt{SECRET\_KEY}, \texttt{FLASK\_DEBUG}. The same image therefore
runs on my laptop, in Compose and against the cloud API without code changes, and no secret is in
the code or the image.

In the container the app is served by \textbf{gunicorn} (1 worker, 4 threads), not by
\cmd{flask run}: Flask's built-in server is single-process, not hardened, and meant for
development. Debug mode is \textbf{off} (\texttt{FLASK\_DEBUG=0}, visible as
\texttt{"debug": false} in \texttt{/health}) because the Werkzeug debugger lets anyone who
triggers an error execute arbitrary Python in the browser, and debug pages leak source code and
configuration. One worker keeps one copy of the model in RAM; threads overlap waiting on the
database and the network.

\shot{s4_4_health}{\texttt{/health} response: database ok, model backend ok, debug false.}
